% \noindent - Input: a phage represented by a score vector in $[0,1]^K$ where $K$ is the number of functional dimensions\\
% - Define the label space $\mathcal{Y}$ (finite set of classes)\\
% - Goal: construct a prediction set $C(x)\subseteq Y$ with the conformal guarantee\\
% - The one-envelope design: build one single envelope on all training phages under their true-label NC transformation, 
% rather than separate envelopes per class

% \noindent - Prediction sets
% - Scores
% - Envelope design

\chapter{Methodology}
\label{sec3: methodology}

In this section, we will develop the mathematical framework used to construct prediction sets with 
guaranteed finite-sample
phage-host classification. Section~\ref{sec3.1: problem formalisation} formalises the prediction task and addresses 
missing data in our score vector. Then, building on the limitations of the CSA method identified in Section~\ref{sec2.4.3: 
csa limitations}, Section~\ref{sec3.2: non-convex envelopes} explains why non-convex envelopes are better suited for 
classification, and Section~\ref{sec3.3: nc transformation} defines the label-dependent nonconformity transformation 
used throughout this report to center valid score distributions near the origin. Section~\ref{sec3.4: envelope 
construction} then presents the three envelope geometries evaluated in this work, Radial, Strip, and Collapsed 1D, and 
derives the calibration scaling factor $\hat{t}$.

\section{Problem Formalization}
\label{sec3.1: problem formalisation}
We will now apply the general conformal prediction setup from Section~\ref{sec2.3: conformal prediction} to our 
host prediction task. Recall from Section~\ref{sec2.1: host prediction problem} that $\mathcal{H} = \{h_1, \ldots, h_N\}$ 
represents the discrete set of $N$ candidate bacterial hosts. Our goal is to build a function $C: \mathcal{X} \to 
\mathcal{P} (\mathcal{H})$ that outputs a set of hosts while maintaining the coverage guarantee in Equation~\eqref{eq: 
core coverage guarantee objective}

In our work, the feature space, $\mathcal{X}$, does not use the raw phage sequences directly. Instead it 
consists of fixed-length score vectors ($K$ dimensional) generated by the protein annotation process introduced in 
Section~\ref{sec2.2.2: plms}. So each score $s_k$ corresponds to a specific functional protein category (e.g., capsid, 
tail, lysin), where $K$ is the total number of categories. A specific phage genome may lack proteins in different 
functional categories, therefore the score vector may have structural missingness:
\begin{equation*}
\mathbf{s} = (s_1, \ldots, s_K)^T \in ([0,1] \cup \{\text{NA}\})^K.
\end{equation*}
\noindent Here, $s_k = \text{NA}$ is the structural absence, which means that for category $k$ no annotated protein 
was identified in the phage sequence, rather than a zero-evidence score or missingness caused by random measurement noise.
This value is missing by design, not due to noise, and hence we need to handle it explicitly rather than simply imputing it.

To retain the mathematical guarantee without imputations, we attach a missingness mask, $\mathbf{m}$, to every 
score vector. We represent each score observation as a paired vector $(\mathbf{s}, \mathbf{m}) \in ([0,1] \cup 
\{\text{NA}\})^K \times \{0,1\}^K$:
\begin{equation*}
\mathbf{m} = (m_1, \ldots, m_K)^T \in \{0,1\}^K, \qquad m_k = \mathbf{1}(s_k \neq \text{NA}).
\end{equation*}

Given a calibration set $D_{\text{cal}} = \{(\mathbf{s}_i, \mathbf{m}_i, H_i)\}_{i=1}^n$ of exchangeable samples,
and an unlabelled test sample $(\mathbf{s}_{n+1}, \mathbf{m}_{n+1})$ drawn exchangeably from an unknown joint probability 
distribution. Our goal is to construct a set-valued prediction mapping $C: ([0,1] \cup \{\text{NA}\})^K \times \{0,1\}^K 
\to \mathcal{P} (\mathcal{H})$ at significance level $\alpha \in (0,1)$ such that:
\begin{equation*}
\label{eq: coverage with missingness}
\mathbb{P}\Big( H_{n+1} \in C(\mathbf{s}_{n+1}, \mathbf{m}_{n+1}) \Big) \ge 1 - \alpha.
\end{equation*}
At the same time, to maximize the preditive efficiency, we want to minimize the expected set size $\mathbb{E}\left[ 
|C(\mathbf{s}, \mathbf{m})| \right]$. The sets $C(\mathbf{s}_{n+1}, \mathbf{m}_{n+1})$ as defined above may be empty 
when no candidate label's transformed score falls within the calibrated envelope. The two prediction tasks in this 
report resolve this case differently \ref{sec3.4: envelope construction}.

The total number of observed functional dimensions for a given phage is given by the $L_1$ norm of the mask, 
$\|\mathbf{m}\|_1 = \sum_{k=1}^K m_k$. Then we have $\mathbf{s}_{\mathbf{m}} \in [0,1]^{\|\mathbf{m}\|_1}$, the sub-vector 
restricted strictly to the observed dimensions where $m_k = 1$. All prediction functions and nonconformity 
scoring operators developed in this report are defined over the product space of pairs $(\mathbf{s}, \mathbf{m})$, rather 
than $\mathbf{s}$ alone. This ensures our nonconformity scores and envelope tests remain mathematically valid for any 
missingness pattern, not just fully observed data, $\mathbf{m} = \mathbf{1}$.

We apply this framework to two host prediction tasks:
\begin{enumerate}[label=(\roman*)]
\item \textbf{Gram-type classification} ($N = 2$): $\mathcal{H} = \{\text{gram-neg}, \text{gram-pos}\}$, evaluated over $K$ 
      functional category dimensions.
\item \textbf{Host-level classification} ($N > 2$): $\mathcal{H}$ represents the full set of bacterial hosts in our 
      dataset. For each host $h \in \mathcal{H}$, the score space is decomposed into host-specific sub-vectors of 
      dimension $K_h$, and an individual envelope $E_h$ is fitted independently on training phages with true host $h$.
\end{enumerate}

In practice, host-level classification is not evaluated over the full host space $\mathcal{H}$ at test time. 
We first obtain a gram-type prediction set $C_{\text{gram}}(\mathbf{s}) \subseteq \{\text{gram-neg}, \text{gram-pos}\}$ 
from the procedure of Section~\ref{sec3.3.1: gram nc mapping}, and use it to restrict the candidate host set to
\begin{equation*}
\hat{\mathcal{H}}(\mathbf{s}) = \bigcup_{g \,\in\, C_{\text{gram}}(\mathbf{s})} \mathcal{H}_g \;\subseteq\; \mathcal{H},
\end{equation*}
where $\mathcal{H}_g \subset \mathcal{H}$ denotes the set of hosts of gram-type $g$. Let $s_h \in \left([0,1] \cup \{NA\}
^{K_h}\right)$ denotes the restriction of $s$ to the $K_h$ score dimensions associated with host $h$. The host-level 
prediction map of this report is therefore the composition
\begin{equation}
\label{eq: host-prediction-set}
C_{\text{host}}(\mathbf{s}, \mathbf{m}) = \big\{ h \in \hat{\mathcal{H}}(\mathbf{s}) \;\big|\; \phi_h(\mathbf{s}_h) 
\in \hat{t}_h \cdot E_h \big\},
\end{equation}
rather than a single-stage test over all of $\mathcal{H}$. This cascade reduces the number of per-host envelope 
evaluations required at test time, at the cost of making $C_{\text{host}}$ dependent on the upstream gram-type 
prediction.
%  we will see the implications of this dependency in Section~\todo{refer}.

The following result confirms that composing the gram type and host level stages does not degrade the 
coverage guarantee of Equation~\eqref{eq: coverage with missingness}.

\begin{proposition}[Cascade Validity]
\label{prop:cascade}
Under exchangeability, if gram type classification is class conditionally valid at level $\alpha$ and each candidate 
host envelope satisfies validity level $\alpha_h$, the cascaded predictor satisfies
\begin{equation}
\label{eq: cascade-bound}
\mathbb{P}\big(H^\star \notin C_{\text{host}}(\mathbf{s})\big) \le \alpha + \max_{h \in \mathcal{H}} \alpha_h.
\end{equation}
\end{proposition}

\begin{proof}
The miscoverage event $\{h^\star \notin C_{\text{host}}(\mathbf{s})\}$ implies either that the true host was pruned by the Gram 
type filter ($h^\star \notin \mathcal{H}_{\text{cand}}(\mathbf{s})$, which requires $g^\star \notin C_{\text{gram}}(\mathbf{s})$) 
or that $h^\star$ was admitted to the candidate pool but rejected by its host envelope ($\tau_{h^\star}(\mathbf{s}) > 
\hat{t}_{h^\star}$). Applying the union bound:
\begin{align*}
\mathbb{P}\big(h^\star \notin C_{\text{host}}(\mathbf{s})\big) &\le \mathbb{P}\big(g^\star \notin C_{\text{gram}}(\mathbf{s})\big) + \mathbb{P}\big(\tau_{h^\star}(\mathbf{s}) > \hat{t}_{h^\star}\big) \\
&= \sum_{g} \mathbb{P}(G=g)\,\mathbb{P}(g^\star \notin C_{\text{gram}} \mid G=g) + \sum_{h} \mathbb{P}(H=h)\,\mathbb{P}(\tau_h > \hat{t}_h \mid H=h) \\
&\le \alpha + \max_{h \in \mathcal{H}} \alpha_h. \qedhere
\end{align*}
\end{proof}

For $\alpha=0.10$, this gives a conservative worst-case coverage lower bound of $1-2\alpha=0.80$; 
Section~\ref{sec5.3: host results} reports the much tighter empirical coverage actually achieved.

